% q0011.tex — Phillips, We Have a Problem (History / Phillips Curve)
\question{
  id        = {0011},
  title     = {Phillips, We Have a Problem},
  source    = {OBECON},
  year      = {2026},
  round     = {Seletiva IEO -- Phase A},
  language  = {English},
  topic     = {Macro},
  subtopic  = {History of Economic Thought / Phillips Curve},
  type      = {objective},
  statement = {%
    1929 and the 1970s were worrisome economic periods for those alive at the
    time, marked respectively by the Great Depression and two oil shocks. In
    both periods, stagnant economic growth was a problem, and unemployment
    rose to sky-high levels.

    Another important factor to consider when looking at an economic crisis is
    inflation. Given your previous knowledge, how was inflation during such
    periods?
  },
  choices   = {%
    \item Prices experienced major deflation in both crises, which decreased
          purchasing power and discouraged spending.
    \item In the Great Depression, prices surged before decreasing rapidly
          due to expansionary monetary policy. In the 1970s, on the other
          hand, prices stayed low for the entire period, except for the rise
          in oil prices.
    \item The Great Depression was a deflationary crisis, while the 1970s saw
          high inflation levels in a phenomenon known as stagflation.
    \item In the Great Depression and in the 1970s, prices suffered
          deflationary pressure, in accordance with the Phillips Curve
          relationship between unemployment and inflation.
    \item Prices surged in both crises, which pushed Central Banks around the
          world to implement contractionary monetary policy.
  },
  answer    = {C},
  solution  = {%
    \begin{itemize}
      \item \textbf{Great Depression (1929):} A severe demand collapse led
            to \textbf{deflation} (prices fell sharply). Unemployment reached
            $\sim$25\% in the US.
      \item \textbf{1970s oil shocks:} Supply-side shocks (OPEC embargo 1973,
            Iranian revolution 1979) pushed costs up, causing simultaneously
            high inflation \emph{and} high unemployment --- a phenomenon
            called \textbf{stagflation}, which broke the traditional
            Phillips Curve trade-off.
    \end{itemize}

    Option (D) is wrong precisely because the 1970s \emph{violated} the
    Phillips Curve: high unemployment coincided with high inflation, not
    deflation.

    Answer: \textbf{(C)}.
  },
}
